\section{Выводы}

Выполнив данную лабораторную работу, я детально изучил и реализовал один из самых сложных строковых алгоритмов — суффиксное дерево алгоритмом Укконена.

\textbf{Об алгоритмике.} Я на практике осознал, как промежуточный алгоритм с асимптотикой $O(N^3)$ математически сводится к $O(N)$ с помощью трех ключевых оптимизаций: использования суффиксных ссылок для телепортации по дереву, введения активной точки (Active Point) и применения ссылок-связей для спуска без посимвольного сравнения.

\textbf{Об оптимизации C++.} Работа с гигантскими структурами данных показала критическую важность профилирования памяти. Изначальная реализация дерева на указателях (\texttt{new Node}) приводила к мгновенным утечкам памяти и падениям системы на больших текстах из-за фрагментации кучи. Перевод графа на индексную адресацию внутри единого \texttt{std::vector} с предварительным резервированием памяти позволил стабилизировать программу.

\textbf{О компромиссах в проектировании.} Проведенный бенчмарк наглядно показал главный инженерный компромисс: Если задача заключается только в поиске минимального сдвига, алгоритм Бута предпочтительнее. Однако суффиксное дерево — это глобальный индекс. Заплатив высокую цену оперативной памятью и временем построения, мы получаем структуру, способную за $O(1)$ или $O(L)$ решать десятки сложнейших задач: от поиска наибольшей общей подстроки до сжатия данных.
\pagebreak